\section{Results}
\label{sec:results}

Our results answer the study's two pre-registered questions in turn, and the two
are not on equal footing. The first is the paper's headline result: does the AIPR
score agree with the human outcome (\S\ref{sec:res-validity}, hypotheses H1--H5)?
The answer is yes, and that validated agreement is what licenses the deployable
triage claim---it stands on its own, regardless of how the second question
resolves. The second asks what the \emph{engineered pipeline} adds over the
underlying model alone---a generic one-paragraph prompt on the same model and PDF
(\S\ref{sec:res-value}, the pre-registered value comparison V1). That is a question
about the \emph{source} of the validated signal, not about whether the signal is
real. Throughout we lead with the large full-mini cohort
($n=\NminiPrimary$) and confirm on the frontier cohort ($n=\NfullPrimary$);
headline metrics for both appear in Table~\ref{tab:headline}.

\subsection{The score agrees with the human outcome}
\label{sec:res-validity}

\begin{table}[t]
  \centering
  \caption{Headline metrics on the primary venue, for the cheap-model full-mini
  configuration (statistical power) and the frontier full configuration
  (production). Brackets are 95\% bootstrap confidence intervals.}
  \label{tab:headline}
  \input{tables/tab_headline}
\end{table}

\paragraph{H1: weak submissions are consistently flagged.}
We report the deployable flag on the production (frontier) model directly. Binning
its scores into quintiles, the lowest-score band is rejected at
\bottomRejectRateFrontier\% (95\% CI \bottomRejectCIFrontier), a
\bottomLiftFrontier$\times$ lift over the \baseRejectRate\% base reject rate, with no
oral papers in that band (\bottomOralRateFrontier\%). We lead with the production
model here, rather than the cheap-model proxy, precisely because the
mini-to-frontier bridge is weakest at the low end where the flag fires
(H5; \S\ref{sec:res-validity}): computing the flag on the frontier score sidesteps
the proxy entirely. The same pattern holds at larger scale on the cheap-model cohort
(sample-relative quintiles; Table~\ref{tab:bands}): the lowest band is rejected at
\bottomRejectRate\% (95\% CI \bottomRejectCI), a \bottomLift$\times$ lift, with the
lift CI \bottomLiftCI{} clearing the pre-registered success rule (lower bound above
one) and oral papers making up \bottomOralRate\% of the band. The reliability
curve (Fig.~\ref{fig:validation}b) shows the empirical reject rate falling
monotonically as the score rises, crossing the base rate near the middle of the
scale. A low score is therefore a high-precision flag that a submission is weak
relative to the venue bar. The bottom-quintile threshold is sample-relative, not
an absolute calibrated cutoff. Because the balanced cohort over-samples the
accept tiers, we also report the flag at the venue's natural prevalence
(\natAcceptRate\% accept; Appendix~\ref{app:robust}): the bottom-quintile reject
precision is \natBottomPrecision\%.

\begin{table}[t]
  \centering
  \caption{Reject rate, lift over base rate, and oral-paper share by AIPR-score
  quintile (full-mini cohort). Band Q1 is the lowest-scoring fifth.}
  \label{tab:bands}
  \input{tables/tab_score_bands}
\end{table}

\begin{figure}[t]
  \centering
  \includegraphics[width=\textwidth]{fig_validation}
  \caption{\textbf{The AIPR score agrees with the human outcome three ways} (full-mini
  cohort). \textbf{(a)} Overall score across the ordered decision tiers (violin,
  box, jittered points), rising monotonically (H3). \textbf{(b)} Reject rate by
  score decile with Wilson 95\% intervals; shaded band = bottom score quintile,
  dashed line = base reject rate---low scores reliably mark rejected work (H1).
  \textbf{(c)} Score vs.\ mean reviewer rating, coloured by tier, with a
  least-squares trend and Spearman $\rho$ (H4). The AUROC scalar (H2) is in
  Table~\ref{tab:headline}; the overlaid ROC curves are Fig.~\ref{fig:naive}a.
  Interactive counterpart (per-point subscores on hover):
  \url{https://aipr.pub/publications\#fig-validation}.}
  \label{fig:validation}
\end{figure}

\paragraph{H2: reject/accept separation.}
The overall score separates rejected from accepted submissions with
AUROC~$=$~\aurocMiniFull on the full-mini cohort, confirmed at \aurocFullFull on the
frontier cohort (Table~\ref{tab:headline}; ROC in Fig.~\ref{fig:naive}a). The
standardized mean gap between accepted and rejected scores is large (Cohen's
$d=\cohensDMini$, Cliff's $\delta=\cliffsDeltaMini$).

\paragraph{H3: monotone quality gradient.}
Mean overall score increases monotonically across the three ordered decision tiers
(Fig.~\ref{fig:validation}a); the Jonckheere--Terpstra trend test rejects the
no-trend null with permutation $\trendPrel$, and the Spearman correlation between
score and tier rank is $\rho=\trendRho$ (95\% CI \trendRhoCI). The gradient is
consistent with the human-review asymmetry: tiers are cleanly ordered, with the
largest separation at the reject boundary.

\paragraph{H4: strong continuous agreement.}
Beyond the discrete decision, the score correlates with the continuous mean
reviewer rating at Spearman $\rho =$~\spearmanRatingMiniFull on the full-mini cohort and
\spearmanRatingFullFull on the frontier cohort (Fig.~\ref{fig:validation}c).
Because the reviewer rating averages over individual reviewer noise, this is the
cleaner test of agreement with human-perceived quality, and the effect is strong
rather than merely significant.

\ifanonymous
\paragraph{The extremes.}
The agreement shows plainly at the ends of the frontier cohort's score range.
AIPR's top-scoring submission was an accepted oral (overall~83), and its
lowest-scoring submission was a reject (overall~58, mean reviewer rating~3.5). The
score places the clearest accept well above the clearest reject, the same ordering
the discrimination and rank agreement above report in aggregate. The disagreement
cases, where the two diverge, are characterized in \S\ref{sec:failure}.
\else
\paragraph{The extremes, named.}
The agreement shows plainly at the ends of the frontier cohort's score range.
AIPR's top-scoring oral was \citet{astabench2026} (overall~83), and its
lowest-scoring reject was \citet{episodicmem2026} (overall~58, mean reviewer
rating~3.5)---both public ICLR~2026 submissions, named here for concreteness. The
score places the clearest accept well above the clearest reject, the same ordering
the discrimination and rank agreement above report in aggregate. The disagreement
cases, where the two diverge, are characterized in \S\ref{sec:failure}.
\fi

\paragraph{H5: the cheap-model proxy is valid.}
On the frontier cohort, the cheap-model full-mini score and the frontier full score
correlate at $\rho =$~\bridgeRhoFull ($n=\bridgeN$; Fig.~\ref{fig:bridge}), clearing
the pre-registered $\rho\ge0.8$ gate. Because a high global correlation can hide
disagreement at the low end where the deployable claim lives, we check the bottom
band directly: within the full-mini bottom quintile the two scores still correlate at
$\rho =$~\bridgeLowRhoFull, and \bridgeOverlap\% of the full-mini bottom quintile is
also in the frontier bottom quintile (Fig.~\ref{fig:bridge}, right). The large-$N$
full-mini results therefore stand in for the production score's behavior, and the
frontier cohort serves as confirmation rather than as the sole evidence.

\subsection{Intelligence is not the bottleneck: what the engineering adds}
\label{sec:res-value}
A validated score raises an obvious question: a frontier model is available to
anyone, so what does the engineered pipeline add over simply asking that model?
This is the study's pre-registered primary value comparison (V1;
\S\ref{sec:methods}, Appendix~\ref{app:naive}): we grade the frontier cohort a
second way with a single one-paragraph prompt that asks the same model, on the
same submitted PDF, for an overall quality score and an accept/reject call. The
prompt mirrors what a researcher would paste into a chat assistant and is
reproduced verbatim in Appendix~\ref{app:naive}, so the baseline is realistic
rather than a strawman.

The first thing the comparison shows is about the model. The bare one-paragraph
prompt, with no rubric and no audit, already separates rejected from accepted
submissions on its own at AUROC \aurocNaiveFull{}---real discrimination against the
human outcome from raw model judgment alone. The full pipeline reaches AUROC
\aurocFullFull{} on the same papers, and the paired difference between them is small
and not statistically significant at conventional levels (\aurocDiffFullNaive,
95\% CI \aurocDiffCI, \aurocDiffPrel; $n=\naiveN$). We therefore do not claim
equivalence or equal performance: the study was powered to detect a larger
pipeline advantage, not to rule in or rule out a small difference over an already
strong naive baseline. Both graders' reject-vs-accept ROC curves clear the
diagonal and overlap closely (Fig.~\ref{fig:naive}a), and the per-tier score
distributions tell the same story (Appendix~\ref{app:robust},
Fig.~\ref{fig:score-dist}).
The score-outcome agreement is therefore not manufactured solely by the engineering:
\emph{with or without our prompt, a frontier model already produces a first-pass
score that substantially tracks the human outcome in this cohort---raw model
intelligence is not the observed bottleneck}, and the reviewer, not the model, keeps
the decision. This also settles the pre-registered V1
prediction, which was directional: we pre-declared that the pipeline would
\emph{out}-discriminate the naive judge, with the AUROC-difference CI excluding zero
(\S\ref{sec:methods}). That criterion is not met---the CI includes zero---so we
report the comparison plainly, as the pre-registration committed us to. The
non-significant \aurocDiffPrel{} result is evidence that the naive baseline is
already strong and that our paired cohort is not powered for small differences; it
is not proof that AIPR and the naive prompt have identical discrimination.
Crucially, this unresolved comparison does not weaken the paper's headline. Both
graders validate against the human outcome---the full pipeline at
AUROC~\aurocFullFull{} and the bare prompt at \aurocNaiveFull{}---so a frontier
model's first-pass score is a useful triage signal whichever of the two produces
it. That high discrimination, not the margin between two already-validated graders,
is the result; the open question V1 leaves is \emph{where} the signal comes from,
not \emph{whether} it is there.

What the engineered pipeline adds is not a higher score but a deployable one.
First, reliability: AIPR grades far more consistently across repeated runs (median
within-paper SD \fullRunSD{} vs.\ \naiveRunSD{} points across repeated gradings of
\runSDnPapers{} papers; Fig.~\ref{fig:naive}b), so the same submission earns nearly
the same verdict rather than a different number each time---the run-to-run noise is
small against the between-tier separation. Second, the same pass that
produces the score produces a structured review---five rubric dimensions, a
grounded citation audit, and verbatim-anchored comments---rather than the bare
number the prompt returns; that review, not the scalar, is what a reviewer acts
on. The scoring rule does several things at once without trading away
discrimination: the validated overall is a by-product of the pass that yields the
review, and the threshold-free AUROC and run-to-run reliability above are where the
discrimination and the value live. The pre-registered AIPR@60 operating point is a
fixed-cutoff diagnostic, not the deployment claim: because the production score is
compressed high and anchored near the venue bar (\S\ref{sec:discussion}), a fixed 60
threshold predicts almost every paper accept and is not venue-calibrated
(\opSixtyFullAcc\% vs.\ \opSixtyNaiveAcc\% balanced accuracy for AIPR and the naive
judge), so we do not read it as an accept/reject classifier. For triage the relevant
error is the opposite conditional---accepted or oral work landing in the low-score
band. On the production frontier score, no accepted
(\lowAcceptCountFrontier/\nAcceptFrontier) or oral
(\lowOralCountFrontier/\nOralFrontier) paper falls in the bottom band; on the larger
full-mini cohort \lowAcceptCountMini/\nAcceptMini accepted and
\lowOralCountMini/\nOralMini oral papers do, the rare cases we read qualitatively in
\S\ref{sec:failure} rather than as a fixed-cutoff accuracy. The value of AIPR is
therefore the processing it wraps around an already-capable model, and the case for
building it does not rest on out-scoring that model.

\begin{figure}[t]
  \centering
  \includegraphics[width=\textwidth]{fig_naive}
  \caption{\textbf{A strong naive baseline, and what the engineering adds}, on the same
  frontier model and submitted PDF (cohort $H$; pre-registered comparison V1).
  \textbf{(a)} Reject-vs-accept ROC for all three graders overlaid---frontier,
  full-mini, and the naive one-paragraph judge---each clearing the diagonal; the
  legend gives AUROC with 95\% CIs and the annotated frontier$-$naive paired
  difference (stratified bootstrap, $n=\naiveN$). AIPR's AUROC is higher, but the
  paired gap is not statistically resolved (\aurocDiffPrel)---neither superiority
  nor equivalence is established. \textbf{(b)} Within-paper score SD over repeated
  runs, one dot per variance-sub-study paper coloured by human decision tier, with
  a median bar---AIPR grades far more consistently. The pipeline's contribution is
  reliability (b) and a grounded, rubric-anchored review in the same pass, not raw
  discrimination (a). Per-tier distributions: Appendix~\ref{app:robust},
  Fig.~\ref{fig:score-dist}; matched-accept-rate operating points:
  Appendix~\ref{app:naive}. Interactive counterpart:
  \url{https://aipr.pub/publications\#fig-naive}.}
  \label{fig:naive}
\end{figure}

\subsection{Secondary analyses}
Per-dimension discrimination (Appendix~\ref{app:robust},
Fig.~\ref{fig:subscore-auroc}) shows novelty and applicability carrying the most
signal, consistent with the weighting of Eq.~\ref{eq:overall}. The citation
dimension carries none: it is pinned at the maximum on \citationPinnedFull\% of
frontier papers (\citationPinnedMini\% on full-mini) and its reject/accept AUROC is
\citationAUROC{} (chance). This is a scoring artifact rather than a property of the
manuscripts: in the original study run, an empty audit result was recorded as a
perfect bibliography. We therefore do not interpret the citation subscore as
validated in this cohort; we diagnose the failure in \S\ref{sec:discussion} and
treat it as a limitation in \S\ref{sec:limits}.
Because citation is the most lightly weighted dimension, the headline is unchanged
when it is dropped from the weighting entirely (AUROC \aurocDropCitationFull{}
vs.\ \aurocFull{} on the frontier cohort; Appendix~\ref{app:robust}). Whether the
audit pass adds signal, run-to-run variance (median within-paper SD
$\approx\runSDmedian$ points), prevalence re-weighting, and weighting robustness
are reported in Appendix~\ref{app:robust}; all point in the same direction. A
second-venue replication (ICLR~2025) is deferred to future work.
